\section[Fargues--Fontaine 曲线 III]{Fargues--Fontaine 曲线 III（11 月 9 日）}

\srcpage{17}

\subsection{回顾}

设 $E$ 是局部域，$\mathcal{O}_E\ni\pi$，$\Fq$ 同前。固定一个完备的非阿基米德代数闭域 $C/\Fq$。

我们已经定义了 Fargues--Fontaine 曲线 $X_{C,E}=Y_{C,E}/\varphi_C^{\mathbb{Z}}$，作为 $E$ 上的 adic 空间。（由于 $E$ 始终固定而 $C$ 会变动，这个记号有时缩写为 $X_C$，甚至 $X$。）这里我们把 $Y_{C,E}$ 定义为 $\Spa W_{\mathcal{O}_E}(\mathcal{O}_C)$ 中 $\pi\neq0$ 且 $[\varpi]\neq0$ 的开子集，其中 $\varpi\in C$ 是伪一致化子。\yizhu{原文用 $\varpi$ 记 $C$ 里的伪一致化子，用 $\pi$ 记 $\mathcal{O}_E$ 的一致化子；$Y_{C,E}$ 的定义同时要求 $\pi\neq0$ 与 $[\varpi]\neq0$。}{A（原文 §4.1 p.17）}

上一讲的定理 3.3 给出如下推论：
\begin{itemize}[leftmargin=2em]
  \item 经典点 $X^{\mathrm{cl}}_{C,E}\subset|X_{C,E}|$ 对应于 $C$ 的 untilt $C^{\sharp}/E$，模去 $\varphi_C^{\mathbb{Z}}$。
  \item 任何连通的 affinoid 开子集 $\Spa A\subset X_{C,E}$ 都具有如下性质：$A$ 是主理想整环（PID），且 $\operatorname{Spm}(A)=X^{\mathrm{cl}}_{C,E}\cap|\Spa A|\subset|X_{C,E}|$。
  \item 对任何经典点 $y\in Y^{\mathrm{cl}}_{C,E}$，存在 $t\in C$ 使 $0<|t|<1$，且 $y=V(\pi-[t])$。（在混合特征情形，$t$ 未必唯一。）
\end{itemize}

\subsection{等晶体}

关于 Fargues--Fontaine 曲线最重要的定理，是其上向量丛的分类。这首先需要一些关于等晶体（isocrystal）的知识。

回忆 $\breveE$ 是 $E$ 的极大无分歧扩域的完备化。具体地，$\breveE=W_{\mathcal{O}_E}(\Fq)[1/\pi]$。

回忆等晶体是一个对 $(V,\varphi)$，其中 $V$ 是有限维 $\breveE$-向量空间，而 $\varphi_V\colon V\xrightarrow{\sim}V$ 是 $\varphi_{\breveE}$-线性的自同构。等晶体构成一个 $E$-线性的 $\otimes$-范畴 $\Isoc_E$。\yizhu{等晶体即 $\breveE$ 上带 $\varphi$-半线性自同构的有限维向量空间；后文的「等斜」指一切子对象的斜率都等于它自身的斜率。}{I（本库 SY02-18 术语表）}

\begin{example}\mbox{}\par
\begin{enumerate}[label=(\arabic*),leftmargin=2.2em]
  \item $\Isoc_E$ 中的单位对象是 $(\breveE,\varphi_{\breveE})$。
  \item 一维对象是 $(\breveE,b\varphi_{\breveE})$，其中 $b\in\breveE^{\times}$。事实上，任何这样的等晶体都同构于 $(\breveE,\pi^n\varphi_{\breveE})$，其中 $n\in\mathbb{Z}$ 唯一。我们称 $n$ 为 $(\breveE,b\varphi_{\breveE})$ 的斜率。这是因为基变换把 $b$ 换成 $a^{-1}b\varphi(a)$，即按「$\varphi$-共轭」改变 $b$；在这个一维情形，可以把它写成 $a^{-1}\varphi(a)b$。现在，容易证明 $\mathcal{O}_{\breveE}^{\times}$ 中的任何元素都能表成 $a^{-1}\varphi(a)$。
\end{enumerate}
\end{example}

又，我们有
\[
\Hom\bigl((\breveE,\pi^{n}\varphi_{\breveE}),(\breveE,\pi^{m}\varphi_{\breveE})\bigr)
=\begin{cases}
E, & m=n,\\
0, & m\neq n.
\end{cases}
\]
事实上，这个 Hom 空间就是 $\breveE^{\varphi=\pi^{m-n}}$。由于 $\varphi$ 保持赋值，只有当 $m-n=0$ 时它才非零。

对任何 $\lambda\in\mathbb{Q}$，写 $\lambda=s/r$，其中 $s,r\in\mathbb{Z}$ 互素、$r>0$。令
\[
(V_\lambda,\varphi_{V_\lambda})
=\left(\breveE^{\,r},
\begin{pmatrix}
0&1&&\\
&0&1&\\
&&\ddots&\ddots\\
\pi^{s}&&&0
\end{pmatrix}\varphi_{\breveE}\right).
\]

\begin{theorem}[Dieudonné--Manin]
$\Isoc_E=\bigoplus_{\lambda\in\mathbb{Q}}\Isoc_E^\lambda$，其中 $\Isoc_E^\lambda$ 是「斜率为 $\lambda$ 的等斜（isoclinic）等晶体」构成的子范畴。范畴 $\Isoc_E^\lambda$ 等价于 $V_\lambda$ 与有限维 $E$-向量空间范畴的张量积。
\end{theorem}

\begin{proof}[证明概要]
对任何非零的 $V=(V,\varphi_V)\in\Isoc_E$，令 $\mu(V)=\dfrac{\deg V}{\rank V}$。
\srcpage{18}
这里 $\operatorname{rank}(V)$ 是底层的 $\breveE$-向量空间的维数，而 $\deg(V)$ 是 $\det V=\wedge^{\operatorname{rank}V}(V)$ 的斜率，后者是秩 $1$ 的等晶体。
秩与次数在短正合列中行为良好，从而给出「斜率的 Harder--Narasimhan 形式化」。事实上，由于上述关于 Hom 的计算，它比通常的情形更强。（通常只能得到单向的 Hom 消没，即一个「半正交分解」。）

我们称 $(V,\varphi_V)$ 是半稳定的（semistable），若对一切子对象 $0\subsetneq(V',\varphi_{V'})\subsetneq(V,\varphi_V)$ 都有 $\mu(V')\le\mu(V)$；若该不等式严格，则称它是稳定的。由形式化的论证，任何对象 $(V,\varphi_V)\in\Isoc_E$ 都有唯一的「Harder--Narasimhan 滤过」，即 $\Isoc_E$ 中一个递减的、分离的、穷尽的滤过
\[
V^{\ge\lambda}\subset V
\]
以 $\lambda\in\mathbb{Q}$ 为指标，使得 $V^{\lambda}:=V^{\ge\lambda}\big/\bigcup_{\lambda'>\lambda}V^{\ge\lambda'}$ 是斜率为 $\lambda$ 的半稳定对象。

现在，在这个特殊情形，斜率函数 $\mu'(V):=-\mu(V)=-\dfrac{\deg V}{\rank V}$ 也给出一个 Harder--Narasimhan 形式化，归根结底仍是因为上述关于 Hom 的计算。这两个滤过典范地互相分裂。因此
\[
\Isoc_E\cong\bigoplus_{\lambda\in\mathbb{Q}}\Isoc_E^\lambda,
\]
其中 $\Isoc_E^\lambda$ 是斜率为 $\lambda$ 的半稳定等晶体构成的全子范畴。

余下要分类 $\Isoc_E^\lambda$。先看 $\lambda=0$。我们想说明 $\Isoc_E^0$ 等价于有限维 $E$-向量空间的范畴，等价由 $W\mapsto(W\otimes_E\breveE,\id\otimes\varphi_{\breveE})$ 给出。容易说明这是全忠实的，因此只需说明一切都由此而来。这归结为说明 $\varphi_V$-不变量足够多：令 $W=V^{\varphi_V=\id}$，便有 $W\otimes_E\breveE\xrightarrow{\sim}V$。

对于 Fargues--Fontaine 曲线上向量丛的分类，论证进行到这一步之前是类似的，此后便分道扬镳。在等晶体的情形，收尾相对容易。想法是说明 $V$ 含有一个 $\varphi_V$-稳定的格（lattice）$L\subset V$，即一个有限自由的 $\mathcal{O}_{\breveE}$-模，它在 $\breveE$ 上生成 $V$。做法是任取一个格，再对 $\varphi_V$ 的各次迭代取它的原像之交。问题在于说明这个交仍是一个格；办法是论证否则它将破坏半稳定性。

现在，在整的情形，可以用 Artin--Schreier 理论说明 $L^{\varphi_L=\id}$ 是 $\mathcal{O}_E$ 上具有正确秩的有限自由模。（先模 $\pi$ 研究，再归纳地往上搭。）

对一般的 $\lambda$，利用：若 $(V,\varphi_V)\in\Isoc_E^\lambda$，则内 Hom $\underline{\Hom}\bigl((V_\lambda,\varphi_{V_\lambda}),(V,\varphi_V)\bigr)\in\Isoc_E^0$，从而化归到 $\lambda=0$。
\end{proof}

\begin{remark}
可以说明 $\End(V_\lambda)=D_\lambda$，即不变量为 $\lambda$ 的中心可除代数。这就完成了 $\Isoc_E$ 的描述。
\end{remark}

\subsection{Fargues--Fontaine 曲线上的向量丛}

注意 $\Fq\hookrightarrow C$。因此 $Y_{C,E}\to\Spa\breveE$，且关于 $\varphi_C\curvearrowright Y_{C,E}$ 与 $\varphi_{\breveE}\curvearrowright\Spa(\breveE)$ 等变。这诱导出一个拉回函子
\[
\Isoc_{\breveE}\longrightarrow
\bigl\{\varphi_C\text{-等变的向量丛}/Y_{C,E}\bigr\}
\xrightarrow[\text{下降}]{\ \cong\ }\mathrm{VB}(X_{C,E}).
\]
我们用 $V\mapsto\mathcal{E}(V)$ 记这个函子。令 $\mathcal{O}_{X_{C,E}}(\lambda):=\mathcal{E}(V_{-\lambda})$。正规化的选取使得 $\mathcal{O}_{X_{C,E}}(1)$ 是「丰沛的」。\yizhu{「丰沛」是代数几何中 ample 的通行译名。$X_{C,E}$ 并非概形，此处的「丰沛」只能作形式类比；定理 4.8 给出它的确切含义。}{A（原文 §4.3 p.18）}

\begin{theorem}[$p$ 进 $E$ 的情形属 Fargues--Fontaine；$E=\Fq((t))$ 的情形属 Hartl--Pink '04；Kedlaya '04]
设 $\mathcal{E}$ 是 $X_{C,E}$ 上的向量丛，则 $\mathcal{E}$ 同构于若干 $\mathcal{O}(\lambda)$ 的直和。等价地，
\srcpage{19}
函子
\[
\Isoc_E\longrightarrow\mathrm{VB}(X_{C,E})
\]
在同构类之间诱导一个双射。

更确切地说，
\begin{enumerate}[label=(\arabic*),leftmargin=2.2em]
  \item 任何 $\mathcal{E}$ 都有 Harder--Narasimhan 滤过 $\mathcal{E}^{\ge\lambda}\subseteq\mathcal{E}$，使得每个
  \[
  \mathcal{E}^{\lambda}:=\mathcal{E}^{\ge\lambda}\Big/\bigcup_{\lambda'>\lambda}\mathcal{E}^{\ge\lambda'}
  \]
  都是斜率为 $\lambda$ 的半稳定对象。
  \item $\Isoc_E^\lambda\xrightarrow{\sim}\mathrm{VB}(X_{C,E})^{\lambda}$，后者是斜率为 $\lambda$ 的半稳定向量丛构成的子范畴。
  \item HN 滤过分裂（但分裂不唯一）。从而上述函子是正合的。
\end{enumerate}
\end{theorem}

定理的第 (1) 部分与所有光滑射影曲线的理论相仿。第 (2) 部分与 $\mathbb{P}^1$ 相似，但与其他曲线不相似。第 (3) 部分与亏格 $g=0,1$ 的曲线相似，与更高亏格的曲线不相似。

\begin{remark}
上述函子远不是范畴等价。$\mathrm{VB}(X_{C,E})$ 中的态射要多得多，例如
\[
H^0(X_{C,E},\mathcal{O}(n))=
\begin{cases}
\text{无穷维 }E\text{-向量空间}, & n>0,\\
E, & n=0,\\
0, & n<0.
\end{cases}
\]
我们有
\[
H^1(X_{C,E},\mathcal{O}(n))=
\begin{cases}
0, & n>0,\\
0, & n=0,\\
\text{无穷维 }E\text{-向量空间}, & n<0.
\end{cases}
\]
这些上同调群具有「什么结构」，我们将在讨论 Banach--Colmez 空间时再回来处理。
\end{remark}

\begin{example}
设 $C^{\sharp}/E$ 是 $C$ 的某个 untilt。这给出一个点 $i\colon\Spa C^{\sharp}\hookrightarrow X_{C,E}$。于是有
\[
0\to I\to\mathcal{O}_{X_{C,E}}\to i_*C^{\sharp}\to0.
\]
可以说明 $I\simeq\mathcal{O}_{X_{C,E}}(-1)$。由注 4.5，有 $H^1(X_{C,E},\mathcal{O}(-1))=C^{\sharp}/E$。这是一个「Banach--Colmez 空间」的例子——它由有限维 $C^{\sharp}$-向量空间与有限维 $E$-向量空间搭成。\yizhu{$C^{\sharp}/E$ 指 Banach--Colmez 空间意义下的商，即 $C^{\sharp}$ 模去 $E$ 的作用，而非域扩张的商。}{A（原文 §I.3 p.19）}
\end{example}

\begin{example}
类似地，把上一例的短正合列用 $\mathcal{O}(1)$ 扭转，得到短正合列
\[
0\to\mathcal{O}_{X_{C,E}}\to\mathcal{O}_{X_{C,E}}(1)\to i_*C^{\sharp}\to0.
\]
取上同调便得到一个扩张
\[
0\to E\to H^0(X_{C,E},\mathcal{O}(1))\to C^{\sharp}\to0.
\]
这应当看作一个「非分裂」扩张，尽管要把它说准确并不容易。
\end{example}

下一个目标是概述分类定理的证明。Hartl--Pink 与 Kedlaya 原来的证明依赖于非常困难的计算。Fargues--Fontaine 的证明是一个显著的简化。然而，它确实用到了 $p$-可除群的非平凡输入。我们将给出一个新的证明，其中没有任何计算，但会大量使用 $p$ 进几何的理论（完满空间（perfectoid space）、菱形与 v-下降）。

我们首先做一些化归，这些化归在所有已知的证明中都是一样的。

\srcpage{20}

\subsection{线丛的分类}

对此，第一步是说明「$\mathcal{O}(1)$ 是丰沛的」。

\begin{theorem}[Kedlaya--Liu '15]
对 $X_{C,E}$ 上的任何向量丛 $\mathcal{E}$，以及一切 $n\gg0$，丛 $\mathcal{E}(n):=\mathcal{E}\otimes_{\mathcal{O}_X}\mathcal{O}_X(n)$ 是整体生成的（globally generated），且 $H^1(X_{C,E},\mathcal{E}(n))=0$。
\end{theorem}

\begin{remark}
由维数原因，对 $i\ge2$ 总有 $H^i(X_{C,E},\mathcal{E})=0$。
\end{remark}

我们略去证明；它是一个相当直接的计算。

\begin{corollary}[GAGA]
令 $P:=\bigoplus_{n\ge0}H^0(X_{C,E},\mathcal{O}(n))$，$X^{\mathrm{alg}}_{C,E}:=\Proj(P)$。则存在局部环化拓扑空间之间的自然映射
\[
X_{C,E}\xrightarrow{\ f\ }X^{\mathrm{alg}}_{C,E},
\]
使得 $f^*\colon\mathrm{VB}(X^{\mathrm{alg}}_{C,E})\xrightarrow{\sim}\mathrm{VB}(X_{C,E})$ 是等价，并且保持上同调。
\end{corollary}

这在什么意义下是一个推论？它来自一个形式化的论证，该论证非常抽象地适用于任何满足定理结论的局部环化空间。

\begin{remark}
概形论意义下的 Fargues--Fontaine 曲线 $X^{\mathrm{alg}}_{C,E}$ 是一个正则概形（regular scheme），是 Noether 的且 Krull 维数为 $1$，局部上是一个主理想整环的谱。这是 Fargues--Fontaine 最初研究的形态。
\end{remark}

\begin{corollary}
任何线丛 $\mathcal{L}\in\Pic(X_{C,E})$ 都同构于 $\mathcal{O}(D)$，其中 $D\in\bigoplus_{x\in X^{\mathrm{cl}}_{C,E}}\mathbb{Z}$ 是某个除子。
\end{corollary}

\begin{proposition}
对任何 $x\in X^{\mathrm{cl}}_{C,E}$，$\mathcal{O}(x):=I_x^{-1}$ 同构于 $\mathcal{O}_{X_{C,E}}(1)$。
\end{proposition}

我们将在下一讲用 Lubin--Tate 理论给出证明。

\begin{corollary}
有 $\mathcal{O}(D)\cong\mathcal{O}(\deg D)$，且 $\Pic(X_{C,E})\cong\mathbb{Z}$。我们规范化该等同，使 $\mathcal{O}(1)\mapsto1$\yizhu{原文作「We normalization the identification」，疑为「We normalize the identification」之误；译文按后者理解。}{A（原文 p.20）}。
\end{corollary}

现在可以定义 $\deg\mathcal{E}$ 为 $\det\mathcal{E}=\wedge^{\operatorname{rank}\mathcal{E}}(\mathcal{E})\in\Pic(X_{C,E})\cong\mathbb{Z}$ 的像。这给出一个 Harder--Narasimhan 形式化，进而得到 Harder--Narasimhan 滤过。这就给出定理 4.4 的第 (1) 部分。

余下要分类斜率为 $0$ 的半稳定丛。容易看出，有限维 $E$-向量空间的范畴全忠实地嵌入 $\mathrm{VB}(X_{C,E})^0$，嵌入由 $W\mapsto W\otimes_E\mathcal{O}_{X_{C,E}}$ 给出。

关键在于：若 $\mathcal{E}/X_{C,E}$ 非零且斜率为 $0$ 半稳定，则 $H^0(X_{C,E},\mathcal{E})\neq0$。

一旦知道这一点，我们就得到从平凡丛出发的一个映射；其余核同样是斜率为 $0$ 的半稳定丛，于是归纳地完成。

想法：我们允许自己（先验地）扩大 $C$。考虑函子 $C'/C\mapsto H^0(X_{C',E},\mathcal{E}_{X_{C',E}})$。这是扩张 $\{C'/C\}$ 上的某个函子。我们希望这个函子可由一个几何对象表示。该对象的 $C'$-值点正是 $H^0(X_{C',E},\mathcal{E}_{X_{C',E}})$。这促使我们考虑比 $C'$ 更一般的测试对象。然而，$X_{C,E}$ 的构造确实依赖于 $C$ 是完满环。所以我们把函子延拓到所有完满 $C$-代数上。这给出（affinoid）完满空间范畴上的一个层。

\begin{example}
固定 $C$ 的一个 untilt $C^{\sharp}/E$。设 $R$ 是完满 $C$-代数。由倾斜等价，存在唯一的完满 untilt $R^{\sharp}/C^{\sharp}$。我们将定义 $X_{R,E}$，使之落入下面的交换图：
\srcpage{21}
\[
\begin{tikzcd}[column sep=1.6em, row sep=2.4em]
\mathrm{Spa}\,R^{\sharp}\arrow[r,"i_{R^{\sharp}}"]\arrow[d]
 & X_{R,E}\arrow[d]
 & & 0\arrow[r] & \mathcal{O}(-1)\arrow[r] & \mathcal{O}\arrow[r]
 & i_{R^{\sharp}*}\mathcal{O}_{R^{\sharp}}\arrow[r] & 0\\
\mathrm{Spa}\,C^{\sharp}\arrow[r,"i_{C^{\sharp}}"]
 & X_{C,E}
 & & 0\arrow[r] & \mathcal{O}(-1)\arrow[r] & \mathcal{O}\arrow[r]
 & i_{C^{\sharp}*}\mathcal{O}_{C^{\sharp}}\arrow[r] & 0
\end{tikzcd}
\]

该正合列表明
\[
H^1(X_{R,E},\mathcal{O}(-1))=R^{\sharp}/E.
\]
于是在等特征情形（此时 $R^{\sharp}=R$），函子 $H^1(X_{C,?},\mathcal{O}(-1))$ 应当是 $\mathbb{A}^1_C/E$。其中 $E\subset\mathbb{A}^1_C$ 通过 untilt 的选取被嵌入为一个闭子集。这是 $\mathbb{A}^1_C$ 对一个 pro-étale 等价关系的商；与代数几何中通常考虑的对象相比，这是病态的。

这将是一般图景：$H^0(X_{?,E},\mathcal{E})$ 与 $H^1(X_{?,E},\mathcal{E})$ 都是完满空间对 pro-étale 等价关系的商。这类对象称为菱形（diamond）。它们与 Artin 的代数空间形式地类似，后者是概形对 étale 等价关系的商。我们将利用菱形的一些抽象性质说明 $H^0(X_{?,E},\mathcal{O}(1))$ 不可能平凡。这保证在某个大的扩张上有截面，而我们可以把它下降。
\end{example}

\begin{example}
在等特征且 $n>0$ 时，$H^0(X_{C,E},\mathcal{O}(n))\cong\mm_{\mathcal{O}_C}^{n}$，它由 $\mathbb{D}^n_C$ 表示。
\end{example}
